\documentclass[10pt,a4paper]{amsart}
\usepackage[margin=19mm]{geometry}
\usepackage{amsmath,amssymb,mathtools}
\usepackage[scr=rsfs,cal=euler]{mathalfa}
\usepackage{xfrac}
\usepackage[unicode,hidelinks,bookmarks=false]{hyperref}
\newcommand{\ie}{i.e.\ }
\newcommand{\cf}{cf.\ }
\newcommand{\resp}{respectively\ }
\newcommand{\Subsection}{\subsection}
\providecommand{\nobreakdash}{\nobreak\hbox{-}}
\setlength{\emergencystretch}{2em}
\multlinegap0pt
\usepackage{bm}
\usepackage{bbm}
\usepackage{dsfont}
\usepackage{xspace}
\usepackage{cancel}
\usepackage{mathtools}
\usepackage{cleveref}
\usepackage{amsthm}
\makeatletter
\@ifundefined{Theorem}{\newtheorem{Theorem}{Theorem}}{}
\@ifundefined{lemma}{\newtheorem{lemma}[Theorem]{Lemma}}{}
\@ifundefined{proposition}{\newtheorem{proposition}[Theorem]{Proposition}}{}
\@ifundefined{theorem}{\newtheorem{theorem}[Theorem]{Theorem}}{}
\makeatother
\makeatletter
\newcommand{\addable}{\alpha}
\newcommand{\cotransitionindices}{\phi}
\newcommand{\cotransitionperms}{\Phi}
\newcommand{\demprod}{\delta}
\newcommand{\divdiff}{N}
\newcommand{\dom}{\mathrm{dom}}
\expandafter\ifx\csname example\endcsname\relax
\newenvironment{example}
{\pushQED{\qed}\renewcommand{\qedsymbol}{$\diamondsuit$}\examplex}
{\popQED\endexamplex}
\fi
\newcommand{\kcotransitionperms}{\overline{\Phi}}
\newcommand{\kdivdiff}{\overline{N}}
\newcommand{\kpipes}{\overline{\sf Pipes}}
\makeatother
\title[Changing Bases with Pipe Dream Combinatorics]{Changing Bases with Pipe Dream Combinatorics}
\author{Anna Weigandt}
\date{Source: arXiv:2506.07306v1 (CC BY 4.0)}
\begin{document}
\maketitle

\noindent\emph{Source and licence.} Changing Bases with Pipe Dream Combinatorics. Version arXiv:2506.07306v1 (CC BY 4.0), distributed under the Creative Commons Attribution 4.0 International licence, as recorded in the arXiv licence metadata for this exact version and shown on the abstract page https://arxiv.org/abs/2506.07306v1. Full rights notice: licenses/arxiv-2506.07306-CC-BY-4.0.txt.\par\smallskip

\noindent\textit{Editorial scope.} Counted unit: the Theorem whose source label is \texttt{thm:pipe\_groth\_to\_schub\_no\_signs} in the article (source0.tex lines 3827--3831, source label \texttt{thm:pipe\_groth\_to\_schub\_no\_signs}), delivered together with the article's own complete original proof of it (source0.tex lines 3833--3875, one \texttt{proof} environment of 43 lines). No mathematics is written, repaired or re-derived: statement and proof are the article's own text line for line. Required context retained uncounted: eq:grothdivdiffbeta (lines 1986--1992); prop:kcotransition (lines 3400--3407); lemma:kcotransitiondescents (lines 3602--3614); lemma:whathappenswithwsianddescents (lines 3634--3645); lemma:cotransitionflippermutation (lines 3681--3687); lemma:kcotransitionrestrictedbijection (lines 3705--3715); eq:grothorddivdiff (lines 3752--3756). Every definition, display and label of those lines is retained unchanged. Omitted from this package and not redistributed: 1--1985, 1993--3399, 3408--3601, 3615--3633, 3646--3680, 3688--3704, 3716--3751, 3757--3826, 3832--3832, 3876--7186. Nothing inside a retained excerpt is omitted or altered: every retained body is delivered exactly as the article's lines are written, so no omission receipt of any kind accompanies this package. Citations: the retained excerpts carry no citation command at all, so no bibliography entry of the article is transcribed and no reference is redistributed. Reference adaptation: none was needed and none was made; every reference inside the delivered unit resolves inside it, so no unresolved reference or citation is produced. Printed slips kept literal: no printed word, symbol, hypothesis or number of the counted statement or of its proof was altered. Layout and class: the article's own class is \texttt{amsart} with packages including amssymb, amsmath, amsfonts, amsthm, graphics, mathrsfs, mathtools, xcolor, geometry, tikz, hyperref, cleveref, thmtools, thm-restate; the wrapper is the lane's pinned \texttt{amsart} editorial wrapper at 10pt on A4, which restates the theorem-like environments the retained text needs and adds the article's own macro definitions that the retained text needs (\texttt{\textbackslash addable}, \texttt{\textbackslash cotransitionindices}, \texttt{\textbackslash cotransitionperms}, \texttt{\textbackslash demprod}, \texttt{\textbackslash divdiff}, \texttt{\textbackslash dom}, \texttt{\textbackslash kcotransitionperms}, \texttt{\textbackslash kdivdiff}, \texttt{\textbackslash kpipes}), with extra wrapper packages (bm, bbm, dsfont, xspace, cancel, mathtools, cleveref). The delivered numbering is the unit's own. Assets: no retained span contains a figure environment, a graphics-inclusion command, a PDF-page inclusion, a table or a TikZ picture; no image file is copied into this package, the package declares no asset dependency and no image file is redistributed with it. Licence: version arXiv:2506.07306v1 of this article is distributed under the Creative Commons Attribution 4.0 International licence, as recorded in the arXiv licence metadata for this exact version and shown on the abstract page https://arxiv.org/abs/2506.07306v1. A full rights notice is delivered at licenses/arxiv-2506.07306-CC-BY-4.0.txt. Characters: every retained excerpt is pure ASCII, so the delivered engine, XeLaTeX with its default Latin Modern text font, reports no missing character.
	\begin{equation}\label{eq:grothdivdiffbeta}
		\kdivdiff^{(\beta)}_i(\mathfrak{G}^{(\beta)}_w) = 
		\begin{cases}
			\mathfrak{G}^{(\beta)}_{ws_i} &\text{if } w > ws_i, \\
			-\beta \mathfrak{G}^{(\beta)}_{w} &\text {if } w < ws_i.
		\end{cases}
	\end{equation}

	\begin{proposition}
		\label{prop:kcotransition}
		Let $w \in S_+$. Suppose $(i,j)$ is an addable cell of $\dom(w)$.
		\begin{enumerate}
			\item $\displaystyle x_i\mathfrak{G}_w = \sum_{u \in \kcotransitionperms_i(w)} (-1)^{\ell(u) - \ell(w) + 1}\mathfrak{G}_u$.
			\item $x_i \displaystyle \mathfrak{S}_w = \sum_{u \in \cotransitionperms_i(w)} \mathfrak{S}_{u}$.
		\end{enumerate}
	\end{proposition}

		\begin{lemma} 
			\label{lemma:kcotransitiondescents}
			Let $w \in S_{n} $ with $w \neq w_0$. Suppose $(i,j)$ is an addable cell of $\dom(w)$ with $j = \addable(w)$.
			The following statements hold:
			\begin{enumerate}
				\item $w < ws_i$.
				\item $i + 1 \in \cotransitionindices_i(w)$.
				\item $\cotransitionindices_i(w)-\{i + 1\} \subseteq \cotransitionindices_{i + 1}(ws_i)$.
				\item $w_{U,i} \in \kcotransitionperms_i(w)$ has a descent at $i$ if and only if $w_{U,i}\geq ws_i$, if and only if $i + 1 \in U$.
				\item $(ws_i)_{U,i+1} \in \kcotransitionperms_{i + 1}(ws_i)$ has a descent at $i$ if and only if $U \subseteq \cotransitionindices_i(w)$. 
				In this case, $(ws_i)_{U,i + 1}=w_{U \cup \{i + 1\},i}$.
			\end{enumerate}
		\end{lemma}

		\begin{lemma}
			\label{lemma:whathappenswithwsianddescents}
			Let $w \in S_n$ with $w \neq w_0$. Suppose that $(i,j)$ is an addable cell of $\dom(w)$ with $j = \addable(w)$. Then:
			\begin{enumerate}
				\item For all $u \in \kcotransitionperms_i(w)$, we have $u \in S_n$.
				\item The cell $(i + 1,j)$ is an addable cell of $\dom(ws_i)$. 
				\item $\{u \in \kcotransitionperms_{i + 1}(ws_i) : u \text{ has a descent at } i\} = \{u \in \kcotransitionperms_{i}(w): u > ws_i\}.$
				\item The map $u \mapsto us_i$ defines a bijection from 
				\[\{v \in \cotransitionperms_{i + 1}(ws_i) : v \text{ has a descent at } i\}\]
				to $\cotransitionperms_{i}(w)-\{ws_i\}$.
			\end{enumerate} 
		\end{lemma}

		\begin{lemma}
			\label{lemma:cotransitionflippermutation}
			Let $w \in S_n$ with $w \neq w_0$. 
			Suppose $(i,j)$ is an addable cell of $\dom(w)$ with $j = \addable(w)$. 
			Let $\mathcal{P} \in \kpipes(w)$ and define $\mathcal{Q} = \mathcal{P} \cup \{(i,j)\}$.
			Then $\demprod(\check{\mathcal{Q}})\square \tau_i = \demprod(\check{\mathcal{P}})$.
		\end{lemma}

		\begin{lemma}
			\label{lemma:kcotransitionrestrictedbijection}
			Let $w \in S_n$ with $w \neq w_0$. 
			Suppose that $(i,j)$ is an addable cell of $\dom(w)$ with $j = \addable(w)$.
			The map $\mathcal{P} \mapsto \mathcal{P} \cup \{(i,j)\}$ defines a bijection from $\{\mathcal{P} \in \kpipes(w) : \check{\mathcal{P}} \text{ is reduced}\}$ to 
			\begin{align}
				\label{eq:bigunionintwoforms}
				\bigsqcup_{\substack{u \in \kcotransitionperms_i(w)\\u\geq ws_i}} &\{\mathcal{Q} \in \kpipes(u):\check {\mathcal{Q}} \text{ is reduced and } \demprod(\check{\mathcal{Q}}) \text{ has a descent at } i\}\\
				& = \bigsqcup_{u \in \kcotransitionperms_i(w)} \{\mathcal{Q} \in \kpipes(u):\check {\mathcal{Q}} \text{ is reduced and } \demprod(\check{\mathcal{Q}}) \text{ has a descent at } i\}. \nonumber
			\end{align}
		\end{lemma}

		\begin{lemma}
			\label{eq:grothorddivdiff}
			Let $w \in S_n$. If $w < ws_i$, then 
			\[\divdiff_i(\mathfrak{G}_{w}) = 0.\]
		\end{lemma}

		\begin{theorem}
			\label{thm:pipe_groth_to_schub_no_signs}
			Let $w \in S_n$. Then
			\[\mathfrak{G}_w^{(1)} = \sum_{\substack{\mathcal{P} \in \kpipes(w)\\ \check{\mathcal{P}} \text{ is reduced}}}\mathfrak{S}_{\demprod(\check{\mathcal{P}})}. \]
		\end{theorem}

		\begin{proof}
			We proceed by reverse induction on Bruhat order. In the base case $w = w_0$, there is a single pipe dream of $w_0$, and it is immediate to check that the statement holds. 
			Now suppose $w \neq w_0$, and assume that the statement holds for all $u \in S_n$ such that $u > w$ in Bruhat order. 
			
			Let $j = \addable(w)$ and $i = w^{-1}(j)$. 
			By \cref{lemma:kcotransitiondescents}, we have $ws_i > w$. By \cref{lemma:whathappenswithwsianddescents}, $ws_i\in S_n$. Thus, $ws_i$ falls within the inductive hypothesis.
			As such,
			\begin{equation}
				\label{eq:grothendieckexpansioninductive}
				\mathfrak{G}_{ws_i}^{(1)} = \sum_{\substack{\mathcal{Q} \in \kpipes(ws_i)\\ \check{\mathcal{Q}} \text{ is reduced}}}\mathfrak{S}_{\demprod(\check{\mathcal{Q}})}.
			\end{equation}
			Because $\kdivdiff_i^{(1)} = \divdiff_i + \divdiff_i x_{i + 1}$, we have
			\begin{align*}
				\mathfrak{G}_w^{(1)}
				&= \kdivdiff_i^{(1)}(\mathfrak{G}_{ws_i}^{(1)}) & \text{(by \cref{eq:grothdivdiffbeta})}\\
				& = \divdiff_i(\mathfrak{G}_{ws_i}^{(1)}) + \divdiff_i (x_{i + 1}\mathfrak{G}_{ws_i}^{(1)})\\
				& = \left(\sum_{\substack{\mathcal{Q} \in \kpipes(ws_i)\\ \check{\mathcal{Q}} \text{ is reduced}}}\divdiff_i(\mathfrak{S}_{\demprod(\check{\mathcal{Q}})})\right) + \divdiff_i( x_{i + 1}\mathfrak{G}_{ws_i}^{(1)}) &\text{(by \cref{eq:grothendieckexpansioninductive})}.
			\end{align*}
			Also,
			\begin{align*}
				\divdiff_i( x_{i + 1}\mathfrak{G}_{ws_i}^{(1)})
				& = \sum_{u \in\kcotransitionperms_{i + 1}(ws_i)}\divdiff_i (\mathfrak{G}_u^{(1)} )&\text{(by \cref{prop:kcotransition} )}\\
				& = \sum_{\substack{u \in\kcotransitionperms_{i + 1}(ws_i)\\ u>us_i}}\divdiff_i (\mathfrak{G}_u^{(1)}) &\text{(by \cref{eq:grothorddivdiff})}\\
				& = \sum_{\substack{u \in\kcotransitionperms_{i}(w)\\u > ws_i}} N_i(\mathfrak{G}_u^{(1)}) & \text{(by Part (3) of \cref{lemma:whathappenswithwsianddescents})}.
			\end{align*}
			By \cref{lemma:whathappenswithwsianddescents}, since $j=\addable(w)$, each $u \in \kcotransitionperms_{i}(w)$ is in $S_n$. 
			In particular, each such $u$ satisfies $u > w$, so we may apply the inductive hypothesis to obtain 
			\begin{align*}
				\divdiff_i( x_{i + 1}\mathfrak{G}_{ws_i}^{(1)})
				&= \sum_{\substack{u \in\kcotransitionperms_{i}(w)\\u > ws_i}} \divdiff_i\left(\sum_{\substack{\mathcal{Q} \in \kpipes(u)\\ 
						\check{\mathcal{Q}} \text{ is reduced}}}\mathfrak{S}_{\demprod(\check{\mathcal{Q}})}\right) \\
				& = \sum_{\substack{u \in\kcotransitionperms_{i}(w)\\u > ws_i}} \sum_{\substack{\mathcal{Q} \in \kpipes(u)\\ 
				\check{\mathcal{Q}} \text{ is reduced}}}\divdiff_i(\mathfrak{S}_{\demprod(\check{\mathcal{Q}})}).
			\end{align*}
			Thus,
			\begin{align*}
				\mathfrak{G}_w^{(1)}& = \sum_{\substack{\mathcal{Q} \in \kpipes(ws_i)\\ \check{\mathcal{Q}} \text{ is reduced}}}\divdiff_i(\mathfrak{S}_{\demprod(\check{\mathcal{Q}})}) + \sum_{\substack{u \in\kcotransitionperms_{i}(w)\\u > ws_i}} \sum_{\substack{\mathcal{Q} \in \kpipes(u)\\ \check{\mathcal{Q}} \text{ is reduced}}}\divdiff_i(\mathfrak{S}_{\demprod(\check{\mathcal{Q}})})\\
				& = \sum_{\substack{u \in\kcotransitionperms_{i}(w)\\u\geq ws_i}} \sum_{\substack{\mathcal{Q} \in \kpipes(u)\\ \check{\mathcal{Q}} \text{ is reduced}}}\divdiff_i(\mathfrak{S}_{\demprod(\check{\mathcal{Q}})}).
			\end{align*}
			Indeed, the only $\mathcal{Q}$'s which contribute nonzero terms to the sum above have the property that $\demprod(\check{\mathcal{Q}})$ has a descent at $i$. 
			Therefore, by \cref{lemma:cotransitionflippermutation} and \cref{lemma:kcotransitionrestrictedbijection},
			\[\mathfrak{G}_w^{(1)} = \sum_{\substack{\mathcal{P} \in \kpipes(w)\\ \check{\mathcal{P}} \text{ is reduced}}}\mathfrak{S}_{\demprod(\check{\mathcal{P}})}. \qedhere\]
		\end{proof}

\end{document}
